\documentclass[10pt,a4paper]{amsart}
\usepackage[margin=19mm]{geometry}
\usepackage{amsmath,amssymb,mathtools}
\usepackage[scr=rsfs,cal=euler]{mathalfa}
\usepackage{xfrac}
\usepackage[unicode,hidelinks,bookmarks=false]{hyperref}
\newcommand{\ie}{i.e.\ }
\newcommand{\cf}{cf.\ }
\newcommand{\resp}{respectively\ }
\newcommand{\Subsection}{\subsection}
\providecommand{\nobreakdash}{\nobreak\hbox{-}}
\setlength{\emergencystretch}{2em}
\multlinegap0pt
\usepackage{bm}
\usepackage{bbm}
\usepackage{dsfont}
\usepackage{xspace}
\usepackage{cancel}
\usepackage{mathtools}
\usepackage{amsthm}
\makeatletter
\@ifundefined{theorem}{\newtheorem{theorem}{Theorem}}{}
\makeatother
\newcommand{\Rdmax}{R^{\max }_{D}}
\newcommand{\Rdstar}{R^{* 2}_{D\sim Z | \mathbf{X}}}
\newcommand{\Rd}{R^2_{D\sim Z | \mathbf{X}}}
\newcommand{\Rymax}{R^{\max }_{Y}}
\newcommand{\Ry}{R^2_{Y\sim Z | D\mathbf{X}}}
\newcommand{\XRVI}{\text{XRVI}}
\newcommand{\df}{df}
\newcommand{\fstarsq}{f^{* 2}_{\alpha, \df-1}}
\newcommand{\fstar}{f^*_{\alpha, \df-1}}
\newcommand{\tstar}{t^*_{\alpha, \df-1}}
\title[On the minimum strength of (unobserved) covariates to overtu]{On the minimum strength of (unobserved) covariates to overturn an insignificant result}
\author{Danielle Tsao, Ronan Perry,, Carlos Cinelli}
\date{Source: arXiv:2408.13901v2 (CC BY 4.0)}
\begin{document}
\maketitle

\noindent\emph{Source and licence.} On the minimum strength of (unobserved) covariates to overturn an insignificant result. Version arXiv:2408.13901v2 (CC BY 4.0), distributed under the Creative Commons Attribution 4.0 International licence, as recorded in the arXiv licence metadata for this exact version and shown on the abstract page https://arxiv.org/abs/2408.13901v2. Full rights notice: licenses/arxiv-2408.13901-CC-BY-4.0.txt.\par\smallskip

\noindent\textit{Editorial scope.} Counted unit: the Theorem whose source label is \texttt{thm:xriv-r} in the article (source0.tex lines 713--735, source label \texttt{thm:xriv-r}), delivered together with the article's own complete original proof of it (source0.tex lines 737--756, one \texttt{proof} environment of 20 lines). No mathematics is written, repaired or re-derived: statement and proof are the article's own text line for line. Required context retained uncounted: thm:XRVI-0 (lines 374--387); thm:XRVI-1 (lines 412--423); xrvi1:d (lines 642--644); maxineq (lines 672--674). Every definition, display and label of those lines is retained unchanged. Omitted from this package and not redistributed: 1--373, 388--411, 424--641, 645--671, 675--712, 736--736, 757--859. Nothing inside a retained excerpt is omitted or altered: every retained body is delivered exactly as the article's lines are written, so no omission receipt of any kind accompanies this package. Citations: the retained excerpts carry no citation command at all, so no bibliography entry of the article is transcribed and no reference is redistributed. Reference adaptation: none was needed and none was made; every reference inside the delivered unit resolves inside it, so no unresolved reference or citation is produced. Printed slips kept literal: no printed word, symbol, hypothesis or number of the counted statement or of its proof was altered. Layout and class: the article's own class is \texttt{imsart} with packages including amsthm, amsmath, amsfonts, amssymb, booktabs, natbib, hyperref; the wrapper is the lane's pinned \texttt{amsart} editorial wrapper at 10pt on A4, which restates the theorem-like environments the retained text needs and adds the article's own macro definitions that the retained text needs (\texttt{\textbackslash Rd}, \texttt{\textbackslash Rdmax}, \texttt{\textbackslash Rdstar}, \texttt{\textbackslash Ry}, \texttt{\textbackslash Rymax}, \texttt{\textbackslash XRVI}, \texttt{\textbackslash df}, \texttt{\textbackslash fstar}, \texttt{\textbackslash fstarsq}, \texttt{\textbackslash tstar}), with extra wrapper packages (bm, bbm, dsfont, xspace, cancel, mathtools). The delivered numbering is the unit's own. Assets: no retained span contains a figure environment, a graphics-inclusion command, a PDF-page inclusion, a table or a TikZ picture; no image file is copied into this package, the package declares no asset dependency and no image file is redistributed with it. Licence: version arXiv:2408.13901v2 of this article is distributed under the Creative Commons Attribution 4.0 International licence, as recorded in the arXiv licence metadata for this exact version and shown on the abstract page https://arxiv.org/abs/2408.13901v2. A full rights notice is delivered at licenses/arxiv-2408.13901-CC-BY-4.0.txt. Characters: every retained excerpt is pure ASCII, so the delivered engine, XeLaTeX with its default Latin Modern text font, reports no missing character.
\begin{theorem}[Closed-form expression for $\XRVI^{0}_{\alpha}$]
\label{thm:XRVI-0}
    Let $f_{r} >0$, then the analytical solution for $\XRVI^{0}_{\alpha}$ is
    \begin{align*}
        \XRVI^{0}_{\alpha}
        &= 
        \begin{cases}
            0, &\text{if } \fstar \leq f_{r},  \\
            \displaystyle 1-\left(\frac{f_{r} }{\fstar}\right)^2,
            &\text{otherwise.} \\
        \end{cases}
    \end{align*}
    If $f_{r} = 0$, there is no value of $\Ry$ capable of overturning an insignificant result.
\end{theorem}

\begin{theorem}[Closed-form expression for $\XRVI^{1}_{\alpha}$]
\label{thm:XRVI-1}
    The analytical solution for $\XRVI^{1}_{\alpha}$ is:
    \begin{align*}
        \XRVI^{1}_{\alpha}
        &= 
        \begin{cases}
            0, &\text{if } \fstar \leq f_{r},  \\
            \dfrac{\fstarsq - f_{r} ^2}{1+\fstarsq}, &\text{otherwise.} \\
        \end{cases}
    \end{align*}
\end{theorem}

     \begin{align}
         \Rdstar = \frac{\Rymax}{f_{r} ^2 + \Rymax}. \label{xrvi1:d}
     \end{align}

\begin{align}
        \tstar \leq t^{\max}_{\boldsymbol{R^2}}. \label{maxineq}
\end{align}

\begin{theorem}[Closed-form expression for $\XRVI^{\Rdmax}_{\alpha}$]
\label{thm:xriv-r}
%
Let $\Rdmax > 0$ or $f_{r} > 0$. Then the analytical expression for $\XRVI^{\Rdmax}_{\alpha}$ is
\begin{align*}
    \XRVI^{\Rdmax}_{\alpha} =
    \begin{cases}
        0, &\text{if} \enspace q \leq 0,\\
        \XRVI^{1}_{\alpha}, &\text{if} \enspace {0 < q \leq \Rdmax}, \\
        \frac{-b + s\sqrt{b^2 - 4ac}}{2a}, &\text{otherwise},
    \end{cases}
\end{align*}
where
\begin{align}
    a &= 1 \label{a}, \\ 
    b &= {\scriptstyle - 2\left[ \frac{\fstarsq - (1-\Rdmax)f_r^2}{\fstarsq + \Rdmax} + \frac{2f_r^2(1-\Rdmax)\Rdmax}{(\fstarsq + \Rdmax)^2} \right]}, \label{b} \\
    c &= { \left[\frac{\fstarsq - (1-\Rdmax)f_r^2}{\fstarsq + \Rdmax}\right]^2,} \label{c} \\ 
    % \XRVI^{1}_{\alpha} &= {\scriptstyle \frac{\fstarsq - f_{r} ^2}{1+\fstarsq},}
    q &= { \frac{\fstarsq - f_r^2}{\fstarsq (1+ f_r^2)}},
\end{align}
and $s\in \{-1, 1\}$ is chosen to yield the valid quadratic root, i.e., $t_{\XRVI,\Rdmax}^{\max}=\tstar$. %When $\Rdmax = 0$ and $0 < f_r <  \fstarsq$, then the quadratic root solution holds.  
If $\Rdmax = 0$ and $f_{r} = 0$, there is no value of $\Ry$ capable of overturning an insignificant result.
\end{theorem}

\begin{proof}[Proof of Theorem~\ref{thm:xriv-r}]
As argued in the proofs of Theorems~\ref{thm:XRVI-0}~and~\ref{thm:XRVI-1}, if $\fstar \leq f_{r} $ (which implies $q\leq 0$), then the minimum strength of $Z$  necessary to attain significance is zero.  If $f_r = 0$ and $\Rdmax = 0$ then by Theorem~\ref{thm:XRVI-0}, we cannot overturn an insignificant result. Otherwise, solving for the minimum value of $\XRVI$ that satisfies (\ref{maxineq}) is equivalent to solving for $\XRVI$ at the equality. 

Now let $f_r>0$. Here we have two cases. Either both $R^2$ values reach the bound or the optimal value of $\Rd$ is an interior point. For the latter case, this means the constraint $\Rd$ is not binding and thus $\XRVI^{\Rdmax}_{\alpha}$ should equal  $\XRVI^{1}_{\alpha}$. Recall that $t$ is concave down with respect to $\Rdmax$; therefore, the optimal value of $\Rd$ is the interior point solution $\Rdstar$, as defined in \eqref{xrvi1:d}, if and only if $\Rdstar \leq \Rdmax$. We can simplify this condition to be of the form,
%
\begin{align*}
     q:=\frac{\fstarsq - f_r^2}{\fstarsq (1+ f_r^2)} \leq \Rdmax.
\end{align*}
It remains to solve for the case where both coordinates reach the bound. Here, the equality in (\ref{maxineq}) simplifies to
    \begin{align}
        {\textstyle \tstar = \frac{f_{r} \sqrt{1-\Rdmax} + \sqrt{\Rdmax\XRVI}}{\sqrt{(1-\XRVI)/(\df-1)}}.} \label{XRVI}
    \end{align}
    
We now solve for $\XRVI$ in (\ref{XRVI}) by squaring both sides and taking the valid root of the quadratic equation. The simplified quadratic form is
\begin{align}
        &{\scriptstyle \XRVI^2 - \left[ \frac{2(\fstarsq - (1-\Rdmax)f_r^2)}{\fstarsq + \Rdmax} + \frac{4f_r^2(1-\Rdmax)\Rdmax}{(\fstarsq + \Rdmax)^2} \right]\XRVI} \nonumber \\
        &{\scriptstyle \quad \quad + \left[\frac{\fstarsq - (1-\Rdmax)f_r^2}{\fstarsq + \Rdmax}\right]^2 = ~0. }
\end{align}
The expressions for $a,b,c$ in (\ref{a})-(\ref{c}) immediately follow from Sridharacharya-Bhaskara's formula for quadratic equations.  Finally, we note that if $f_r=0$ and $\Rdmax > 0$ then the $\XRVI^{\Rdmax}$ solution is valid. 
\end{proof}

\end{document}
